\documentclass[11pt,a4paper]{article}

% === PACKAGES ===
\usepackage[utf8]{inputenc}
\usepackage[margin=1in]{geometry}
\usepackage{amsmath,amssymb,amsthm}
\usepackage{enumitem}
\usepackage{graphicx}
\usepackage{hyperref}
\usepackage{fancyhdr}
\usepackage{xcolor}

% === DOCUMENT INFO (CUSTOMIZE THESE) ===
\newcommand{\coursename}{[COURSE NAME]}
\newcommand{\coursenumber}{[COURSE NUMBER]}
\newcommand{\lecturenum}{[LECTURE NUMBER]}
\newcommand{\lecturetopic}{[LECTURE TOPIC]}
\newcommand{\semester}{[SEMESTER]}

% === HEADER/FOOTER ===
\pagestyle{fancy}
\fancyhf{}
\lhead{\coursename}
\chead{Practice Problems}
\rhead{Lecture \lecturenum}
\rfoot{Page \thepage}

% === DIFFICULTY STARS ===
\newcommand{\easy}{{\color{green!70!black}$\bigstar$}}
\newcommand{\medium}{{\color{orange!80!black}$\bigstar\bigstar$}}
\newcommand{\hard}{{\color{red!70!black}$\bigstar\bigstar\bigstar$}}

% === COMMON MATH OPERATORS ===
\DeclareMathOperator{\E}{\mathbb{E}}
\DeclareMathOperator{\Var}{Var}
\DeclareMathOperator{\Cov}{Cov}
\newcommand{\R}{\mathbb{R}}
\newcommand{\N}{\mathbb{N}}
\newcommand{\Z}{\mathbb{Z}}

% === TITLE ===
\title{\textbf{Practice Problems} \\ \coursename{} (\coursenumber) \\ Lecture \lecturenum: \lecturetopic}
\author{}
\date{\semester}

\begin{document}

\maketitle

% ============================================================
% READING ASSIGNMENTS
% ============================================================
\section*{Reading Assignments}

\textbf{Primary:}
\begin{itemize}
  \item [Textbook 1] Chapter X, \S X.1--X.3 ([topic]).\\
        \textit{Focus on:} Definition X.1, Theorem X.2.
  \item [Textbook 2] Chapter Y, \S Y.1--Y.2 ([alternative treatment]).
\end{itemize}

\textbf{Review Questions:}
\begin{itemize}
  \item [Textbook 1] Ch. X end-of-chapter: Exercises 1--5, 8, 12.
  \item [Textbook 2] Ch. Y: Problems 2, 4.
\end{itemize}

\textbf{Optional (Advanced):}
\begin{itemize}
  \item [Paper/Resource]: [Title], pp. XX--YY (for deeper treatment of [topic]).
\end{itemize}

\newpage

% ============================================================
% PROBLEMS
% ============================================================
\section*{Problems}

% --- EASY PROBLEMS ---

\subsection*{Foundational}

\textbf{Problem 1} \easy{} \textsc{[Topic]} \hfill \textit{Source: [Textbook] Ch. X, Ex. 3}

See [Textbook], Chapter X, Exercise 3.

\bigskip

\textbf{Problem 2} \easy{} \textsc{[Topic]} \hfill \textit{Source: Course PS1, Problem 2}

[Copy full problem statement from course problem set]

\bigskip

% --- MEDIUM PROBLEMS ---

\subsection*{Intermediate}

\textbf{Problem 3} \medium{} \textsc{[Topic]} \hfill \textit{Source: [Textbook] Ch. X, Ex. 8}

See [Textbook], Chapter X, Exercise 8.

\bigskip

\textbf{Problem 4} \medium{} \textsc{[Topic]} \hfill \textit{Source: Course PS2, Problem 4}

[Copy full problem statement from course problem set]

\bigskip

\textbf{Problem 5} \medium{} \textsc{[Topic]} \hfill \textit{Source: Original}

[Full problem statement for original problem]

\begin{enumerate}[(a)]
  \item [Part a]
  \item [Part b]
  \item [Part c]
\end{enumerate}

\bigskip

% --- HARD PROBLEMS ---

\subsection*{Challenging}

\textbf{Problem 6} \hard{} \textsc{[Topic]} \hfill \textit{Source: Course Midterm, Problem 3}

[Copy full problem statement from course exam]

\bigskip

\textbf{Problem 7} \hard{} \textsc{[Topic]} \hfill \textit{Source: [Textbook] Ch. X, Ex. 15}

See [Textbook], Chapter X, Exercise 15.

\bigskip

\textbf{Problem 8} \hard{} \textsc{[Topic]} \hfill \textit{Source: Original}

[Full problem statement for challenging original problem]

\bigskip

% ============================================================
% DIFFICULTY KEY
% ============================================================
\vfill
\hrule
\smallskip
\textbf{Difficulty Key:} 
\easy{} Easy (direct application) \quad
\medium{} Medium (multi-step reasoning) \quad
\hard{} Hard (challenging/synthesis)

\end{document}
